\documentclass[letterpaper, 10 pt, conference]{ieeeconf}
\IEEEoverridecommandlockouts
\overrideIEEEmargins

\usepackage{amsmath,amssymb,amsfonts}
\usepackage{graphicx}
\usepackage{cite}

\title{\bf The Rigel Framework: Recoverability Geometry on the Normalized Simplex}
\author{[Authors anonymized for submission]}

\begin{document}
\maketitle
\thispagestyle{empty}
\pagestyle{empty}

%%%%%%%%%%%%%%%%%%%%%%%%%%%%%%%%%%%%%%%%%%%%%%%%%%%%%%%%%%%%%%%%%%%%%%%%%%%%%%%%
\begin{abstract}
We introduce the Rigel framework for analyzing recoverability in systems governed by capacity-constrained autonomy. Building on the GCAT state space, we normalize system states onto the simplex $\Delta^3$ and define a geometric notion of recoverability relative to collapse modes. The Rigel equilibrium serves as a reference point maximizing distance from collapse boundaries. We define the Rigel number as a ratio capturing distance to collapse and distance from equilibrium, and establish its monotonic properties. We show that recoverability is independent of static admissibility and characterize their relationship. Under projected linear dynamics, we establish asymptotic convergence toward the Rigel equilibrium. The framework provides a geometric basis for understanding recoverability prior to enforcement.
\end{abstract}

%%%%%%%%%%%%%%%%%%%%%%%%%%%%%%%%%%%%%%%%%%%%%%%%%%%%%%%%%%%%%%%%%%%%%%%%%%%%%%%%
\section{Introduction}

While admissibility defines whether a system state is sustainable under capacity constraints, it does not characterize whether a system can recover once those constraints are violated. Systems may enter inadmissible states yet remain recoverable, or remain admissible while approaching irreversible collapse.

This paper introduces a geometric framework for recoverability independent of static admissibility. The key idea is to represent system states on a normalized simplex, where proximity to collapse modes and distance from equilibrium define recoverability.

The Rigel framework complements the GCAT formulation by introducing a second axis of analysis: recoverability as a function of geometric position within the simplex.

%%%%%%%%%%%%%%%%%%%%%%%%%%%%%%%%%%%%%%%%%%%%%%%%%%%%%%%%%%%%%%%%%%%%%%%%%%%%%%%%
\section{Normalized Geometry}

\begin{definition}[Normalized State]
For $\mathbf{x} \in \mathcal{X} \setminus \{\mathbf{0}\}$:
\begin{equation}
\tilde{\mathbf{x}} = \frac{\mathbf{x}}{\|\mathbf{x}\|_1} \in \Delta^3
\end{equation}
\end{definition}

The simplex:
\begin{equation}
\Delta^3 = \left\{ \tilde{\mathbf{x}} \in \mathbb{R}^4_+ : \sum_{i=1}^4 \tilde{x}_i = 1 \right\}
\end{equation}

represents relative allocation across governance, control, autonomy, and trust.

%%%%%%%%%%%%%%%%%%%%%%%%%%%%%%%%%%%%%%%%%%%%%%%%%%%%%%%%%%%%%%%%%%%%%%%%%%%%%%%%
\section{Rigel Equilibrium}

\begin{definition}[Rigel Point]
\begin{equation}
\mathbf{R} = \left(\frac{1}{4}, \frac{1}{4}, \frac{1}{4}, \frac{1}{4}\right)
\end{equation}
\end{definition}

\begin{proposition}
$\mathbf{R}$ maximizes the minimum coordinate over $\Delta^3$:
\begin{equation}
\mathbf{R} = \arg\max_{\tilde{\mathbf{x}} \in \Delta^3} \min_i \tilde{x}_i
\end{equation}
\end{proposition}

\textit{Argument.}
By symmetry and convexity of the simplex, the centroid equalizes all coordinates, maximizing the minimum.

%%%%%%%%%%%%%%%%%%%%%%%%%%%%%%%%%%%%%%%%%%%%%%%%%%%%%%%%%%%%%%%%%%%%%%%%%%%%%%%%
\section{Collapse Geometry}

\begin{definition}[Collapse Faces]
\begin{align}
C_1 &= \{\tilde{g}=0\}, \quad
C_2 = \{\tilde{c}=0\}, \\
C_3 &= \{\tilde{a}=0\}, \quad
C_4 = \{\tilde{t}=0\}
\end{align}
\end{definition}

These define the boundaries where one system component collapses.

\begin{definition}[Face Distance]
\begin{equation}
d_F(\tilde{\mathbf{x}}) = \min_i \tilde{x}_i
\end{equation}
\end{definition}

%%%%%%%%%%%%%%%%%%%%%%%%%%%%%%%%%%%%%%%%%%%%%%%%%%%%%%%%%%%%%%%%%%%%%%%%%%%%%%%%
\section{Distance to Equilibrium}

\begin{definition}[Simplex Distance]
\begin{equation}
d_\Delta(\tilde{\mathbf{x}}, \mathbf{R}) = \|\tilde{\mathbf{x}} - \mathbf{R}\|_2
\end{equation}
\end{definition}

This measures deviation from balanced allocation.

%%%%%%%%%%%%%%%%%%%%%%%%%%%%%%%%%%%%%%%%%%%%%%%%%%%%%%%%%%%%%%%%%%%%%%%%%%%%%%%%
\section{The Rigel Number}

\begin{definition}[Rigel Number]
\begin{equation}
R_i(\tilde{\mathbf{x}}) = \frac{d_F(\tilde{\mathbf{x}})}{d_F(\tilde{\mathbf{x}}) + d_\Delta(\tilde{\mathbf{x}}, \mathbf{R}) + \delta}
\end{equation}
\end{definition}

with $\delta > 0$.

\begin{proposition}[Bounds]
\begin{equation}
0 \le R_i(\tilde{\mathbf{x}}) \le 1
\end{equation}
\end{proposition}

\begin{proposition}[Monotonicity]
$R_i$ is increasing in $d_F$ and decreasing in $d_\Delta$.
\end{proposition}

\textit{Argument.}
Both follow directly from partial derivatives of the ratio.

\begin{proposition}[Extrema]
\begin{itemize}
\item $R_i \to 0$ as $\tilde{\mathbf{x}} \to \mathcal{C}$
\item $R_i$ maximized at $\mathbf{R}$
\end{itemize}
\end{proposition}

%%%%%%%%%%%%%%%%%%%%%%%%%%%%%%%%%%%%%%%%%%%%%%%%%%%%%%%%%%%%%%%%%%%%%%%%%%%%%%%%
\section{Recoverability vs Admissibility}

\begin{definition}[Recoverable Region]
\begin{equation}
\mathcal{R} = \{ \tilde{\mathbf{x}} : R_i(\tilde{\mathbf{x}}) > \epsilon \}
\end{equation}
\end{definition}

\begin{proposition}
Recoverability and admissibility are independent:
\begin{equation}
\mathcal{A} \not\subseteq \mathcal{R}, \quad \mathcal{R} \not\subseteq \mathcal{A}
\end{equation}
\end{proposition}

\textit{Interpretation.}
Systems may be admissible but fragile, or inadmissible but recoverable.

%%%%%%%%%%%%%%%%%%%%%%%%%%%%%%%%%%%%%%%%%%%%%%%%%%%%%%%%%%%%%%%%%%%%%%%%%%%%%%%%
\section{Projected Dynamics}

We consider dynamics restricted to $\Delta^3$:
\begin{equation}
\dot{\tilde{\mathbf{x}}} = \Pi_{T\Delta^3}(-K(\tilde{\mathbf{x}} - \mathbf{R}))
\end{equation}

where $\Pi_{T\Delta^3}$ denotes projection onto the tangent space:
\begin{equation}
T\Delta^3 = \left\{ v : \sum_i v_i = 0 \right\}
\end{equation}

%%%%%%%%%%%%%%%%%%%%%%%%%%%%%%%%%%%%%%%%%%%%%%%%%%%%%%%%%%%%%%%%%%%%%%%%%%%%%%%%
\section{Stability Result}

\begin{theorem}
Under the projected dynamics, $\mathbf{R}$ is asymptotically stable.
\end{theorem}

\textit{Argument.}
The system reduces to a linear contraction in the tangent space. Projection preserves feasibility, and $\mathbf{R}$ is the unique equilibrium.

%%%%%%%%%%%%%%%%%%%%%%%%%%%%%%%%%%%%%%%%%%%%%%%%%%%%%%%%%%%%%%%%%%%%%%%%%%%%%%%%
\section{Benchmarks}

\begin{table}[h]
\centering
\caption{Rigel Numbers}
\begin{tabular}{lcccccc}
System & $\tilde{\mathbf{x}}$ & $d_F$ & $d_\Delta$ & $R_i$ & Label \\
\hline
Drone & (0.222,0.333,0.278,0.167) & 0.167 & 0.124 & 0.49 & Marginal \\
Human & (0.270,0.216,0.189,0.324) & 0.189 & 0.104 & 0.55 & Moderate \\
Consensus & (0.400,0.100,0.075,0.425) & 0.075 & 0.326 & 0.17 & Rigid \\
\end{tabular}
\end{table}

These illustrate that recoverability is not aligned with admissibility.

%%%%%%%%%%%%%%%%%%%%%%%%%%%%%%%%%%%%%%%%%%%%%%%%%%%%%%%%%%%%%%%%%%%%%%%%%%%%%%%%
\section{Discussion}

The Rigel framework introduces recoverability as a geometric property distinct from admissibility. By operating on the simplex, it captures structural resilience independent of absolute scale.

%%%%%%%%%%%%%%%%%%%%%%%%%%%%%%%%%%%%%%%%%%%%%%%%%%%%%%%%%%%%%%%%%%%%%%%%%%%%%%%%
\section{Conclusion}

We presented a geometric formulation of recoverability on the GCAT simplex. The Rigel number captures proximity to collapse and distance from equilibrium. This provides a foundation for integrating recoverability into dynamic and enforcement frameworks in subsequent work.

\end{document}
